% year: תשפ"ג
% type: מבחן
% semester: א
% moed: ב
% number: 5ב
% points: 10
% categories: ode
% source: exams/מבחנים/תשפג/מועד ב תשפג.docx

\begin{question}
פתרו את המשוואה $e^x\,dx-y\,dy=0$ עם תנאי התחלה $y(0)=1$.
\end{question}

\begin{hint}
זו משוואה פרידה: העבירו אגף ובצעו אינטגרציה לכל צד בנפרד.
\end{hint}

\begin{solution}
המשוואה פרידה: $y\,dy=e^x\,dx$ (כלומר $y\,y'=e^x$). נבצע אינטגרציה לשני האגפים:
\[ \int y\,dy=\int e^x\,dx\ \Longrightarrow\ \frac{y^2}{2}=e^x+C. \]
מתנאי ההתחלה $y(0)=1$: $\frac12=1+C$, ולכן $C=-\frac12$, ו-
\[ y^2=2e^x-1. \]
מכיוון ש-$y(0)=1>0$ ופתרון רציף אינו יכול לעבור לערכים שליליים בלי לעבור ב-$0$, נבחר את הענף החיובי:
\[ y(x)=\sqrt{2e^x-1}, \]
המוגדר כאשר $2e^x-1>0$, כלומר $x>-\ln 2$.
בדיקה: $y'=\frac{e^x}{\sqrt{2e^x-1}}$, ולכן $y\,y'=e^x$ ✓, ו-$y(0)=\sqrt{1}=1$ ✓.
\end{solution}
